\section{Introducción}\label{sec:Introducción}

\begin{frame}
    \frametitle{Introducción}
    \begin{tcolorbox}[colframe=color1]
        \begin{center}
        ¿Qué constituye un sitio web?
        \end{center}
    \end{tcolorbox}
    \begin{itemize}
        \item Una interfaz gráfica
        \item ¿Nada más?
    \end{itemize}
\end{frame}

\begin{frame}
\frametitle{Introducción}
    \begin{tcolorbox}[colframe=color1]
        \begin{center}
            Front-end
        \end{center}
    \end{tcolorbox}
    
    Cuando nos referimos al desarrollo ``front-end'', estamos hablando
    de el trabajo de programar la interfaz que van a ver los usuarios 
    cuando entran al sitio. Eso podría también podría incluir hacer 
    el diseño, pero no necesariamente. 
    
    Lo que se arma en esta parte podría entenderse como un ``cliente'', un 
    programa que corre en la máquina del usuario y que se comunica 
    con el servidor solo cuando sea necesario a través de cierta ``API''.

    Este taller va a ser una introducción a front-end.

\end{frame}

\begin{frame}
    \frametitle{Introducción}
    \begin{tcolorbox}[colframe=color1]
        \begin{center}
            ¿Servidor? ¿API? Back-end
        \end{center}
    \end{tcolorbox}

    Cuando un sitio se vuelve relativamente complejo, no puede correr 
    o almacenarse entero en la máquina del usuario. No solo por un tema 
    de recursos, sino que también por seguridad. Por ejemplo, en un sitio 
    con autenticación, sería super inseguro y pesado almacenar todas las 
    cuentas y contraseñas en las máquinas de los usuarios. Aparte, el 
    sitio mismo tiene que estar subido a algún lado para que la gente pueda 
    acceder.

    Todo ese trabajo entra en el campo del back-end. Servidores, bases de 
    datos, redes, entre otras cosas, son todas parte de nuestro trabajo 
    como desarrolladores de software. Dependiendo del tamaño y objetivo del 
    sitio, suele ser la mayor parte del desarrollo.

\end{frame}

\begin{frame}
    \frametitle{Introducción}

    En este taller no nos vamos a meter con back-end, pero si les interesa 
    quédense atentes al taller de APIs! También la mayoría de los temas 
    relacionados se ven con muchísima profundidad a lo largo de la carrera 
    de computación.

\end{frame}

\begin{frame}
    \frametitle{Volviendo al front...}

    ¿Y qué significa ``desarrollar front''? ¿Cómo se programan estos clientes? 
    ¿Qué lenguaje y herramientas usamos?

    Los sitios web corren en navegadores web. Estos son Chrome, Firefox, etc. 
    Estos navegadores entienden tres lenguajes:

    \begin{itemize}
        \item \textbf{HTML}: Define la estructura del sitio, su contenido$^*$. 
              En principio esta es la parte más importante, todo lo demás 
              depende e interactúa con esta información.
        \item \textbf{CSS}: Define como se ve el sitio, su estilo. Tiene muchísimas 
              opciones, lo cual lo hace bastante difícil de dominar, pero es 
              necesario para hacer un sitio que no sea un texto negro sobre fondo 
              blanco.
        \item \textbf{JavaScript}: A diferencia de los anteriores, este es un 
              lenguaje de programación. Como tal, se ocupa de la lógica e 
              interactividad del sitio. Es la parte más compleja de muchos 
              de los sitios modernos y eleva las posibilidades a otro nivel.
    \end{itemize}

\end{frame}

\begin{frame}
    \frametitle{Volviendo al front...}

    En este taller vamos a ver solo HTML y CSS. Por separado ya son super 
    potentes y explicar JavaScript implicaría enseñar a programar, para 
    lo cual existe la carrera.

    También, como veremos más adelante, muchas veces no se programa en 
    estos lenguajes de forma directa. Pero nos estamos adelantando, 
    arranquemos desde el principio.

\end{frame}
